\section{Introduction}

Cortical synapses must change in ways that improve behaviour, using only information available
locally. Learning rules that respect this constraint, such as three-factor rules with eligibility
traces and broadcast error signals \citep{gerstner2018,murray2019rflo,bellec2020eprop}, can solve
tasks that otherwise need back-propagation through time (BPTT). They differ from exact gradients in
how credit is assigned \citep{lillicrap2020,richards2019}. If brains learn with such rules, the
difference should leave traces in what a circuit ends up representing. Representations of
networks trained with different rules can indeed be told apart from their observables
\citep{nayebi2020}. Whether these signatures appear in biologically constrained circuits, and
whether they move a model towards or away from real cortex, is unknown.

Mouse primary visual cortex (V1) offers the data needed to ask this. There are fitted
single-neuron models for hundreds of recorded cells \citep{yeung2025gfr}, cell-type-specific
connectivity from electron microscopy and synaptic physiology \citep{campagnola2022}, and
large-scale Neuropixels recordings of responses to the same natural images and gratings
\citep{siegle2021,devries2020}. Building on these, \citet{ito2026} trained a large V1 model with BPTT
to match in-vivo activity statistics. That work leaves open the question we address: does the
learning algorithm itself shape the population geometry, and is a cell-local rule closer to or
further from mouse V1 than the exact gradient?

Our earlier version of this question used abstract working-memory tasks. Its conclusions were
dominated by the task and by the cells' slow intrinsic filters, and its inputs had no counterpart
in any recording. Here we rebuilt the setup around the data. A V1 microcircuit of 1,000
patch-clamp-fitted cells with measured connectivity receives the same visual stimuli as the
recorded mice, through a fixed LGN front end. The same circuit is trained with exact BPTT or with a
cell-local e-prop rule, with or without the cell's own spike-history term. Both use the same
optimizer, learning rates and budget, under two different objectives: matching in-vivo response
statistics, and a Visual-Behavior-style change-detection task. We then compare population
responses to 118 natural scenes that were never used in training with those of mouse V1 units. We
use representational similarity, CKA, dimensionality, sparseness and decoding, correct for the
reliability of both model and data, and relate geometry to task performance along training. The
protocol was locked before the held-out cohort of fitted cells was used.

The main results are:
\begin{itemize}
\item The reliability-corrected similarity to mouse V1 was the same for every rule and for untrained
networks. The shared geometry is set by the circuit's constraints, not by its learning rule.
\item With in-vivo statistics as the objective, the local rule fitted them more closely than BPTT,
but only by moving into a regime with very large weights. There natural-image dimensionality,
decodability and sparseness departed from the mouse. At matched loss, the two rules' geometry was
the same.
\item A behavioral objective changed the natural-scene geometry little under any rule.
\end{itemize}
